%-----------------------------------------------------------------------
\subsection{Instance 6: Model Selection Under the Rashomon Effect}
\label{sec:instance-selection}
%-----------------------------------------------------------------------

Model selection---choosing a single best model from a trained set---is perhaps
the most commonplace decision in applied machine learning.  Yet when many models
achieve near-identical performance, any selection procedure is navigating a
\emph{Rashomon set}: a collection of models that are, by the available evidence,
equally defensible choices.  We show that this setting is an instance of
Theorem~\ref{thm:universal}, and that the resulting instability is not a
failure of any particular selection criterion but a mathematical inevitability.

\paragraph{Explanation system.}
We instantiate the abstract framework as follows.

\begin{definition}[Model Selection Explanation System]
\label{def:ms-system}
The \emph{model selection explanation system} is the tuple
$\mathcal{S}_{\mathrm{sel}} = (\Params, \sH, Y, \mathsf{obs}, \mathsf{exp}, \incomp)$ where:
\begin{itemize}
  \item $\Params = \mathbf{\Theta}$ is the space of \emph{model configurations}
    (hyperparameter settings, architecture choices, training seeds, or any
    parameterization of the model class under comparison);

  \item $\sH$ is the \emph{hypothesis space} of trained models; each
    $h \in \sH$ corresponds to a configuration $\theta(h) \in \Params$ via
    the attribution map $\theta: \sH \to \Params$;

  \item $Y$ is the space of \emph{performance metrics} (validation loss,
    accuracy, AUC, or any scalar summary of model quality on a held-out
    evaluation set); formally, $Y$ is equipped with a strict total order
    $>$ representing ``performs better than'';

  \item $\mathsf{obs}: \sH \to Y$ maps each hypothesis to its
    \emph{observed performance metric} on the evaluation set;

  \item $\mathsf{exp}: \sH \to Y$ is the \emph{selection rule}---the
    criterion used to rank or choose among models; and

  \item $\incomp\;\subseteq Y \times Y$ is defined by $(y, y') \in \incomp$
    whenever $y \neq y'$ and the two metrics imply opposite selection
    decisions (i.e., the model recommended under $y$ is strictly
    disfavored under $y'$).
\end{itemize}
\end{definition}

\paragraph{The Rashomon property for model selection.}
The central empirical finding motivating the Rashomon effect is that, across
many supervised learning tasks, large sets of models achieve near-optimal
performance while differing substantially in structure, predictions, and
feature usage \citep{breiman2001random,fisher2019models,rudin2024amazing}.
Formally, define the \emph{$\varepsilon$-Rashomon set} as
\[
  \mathcal{R}_\varepsilon = \bigl\{ h \in \sH : \mathsf{obs}(h) \geq \mathsf{obs}(h^*) - \varepsilon \bigr\},
\]
where $h^*$ is the empirical risk minimizer and $\varepsilon > 0$ is a
tolerance.  The Rashomon property holds for $\mathcal{S}_{\mathrm{sel}}$
whenever $\mathcal{R}_\varepsilon$ contains pairs of models that are
(i)~indistinguishable in performance ($\theta(h) = \theta(h')$ in the sense
that they share the same configuration class) yet
(ii)~produce incompatible selection metrics on different evaluation samples.

\begin{proposition}[Rashomon Property for Model Selection]
\label{prop:ms-rashomon}
If the $\varepsilon$-Rashomon set $\mathcal{R}_\varepsilon$ contains
two models $h, h' \in \sH$ with $\theta(h) = \theta(h')$ such that
$\mathsf{obs}(h) \neq \mathsf{obs}(h')$, then $\mathcal{S}_{\mathrm{sel}}$
satisfies the Rashomon property (Definition~\ref{def:rashomon}).
\end{proposition}

\noindent
This is immediate from the definition: $\mathsf{obs}(h)$ and
$\mathsf{obs}(h')$ are incompatible (one favors $h$, the other $h'$), and
both are witnessed by models sharing the same underlying configuration.
The condition is generically satisfied: \citet{fisher2019models} show
that the Rashomon set has positive measure for any model class with more
free parameters than identifiable constraints, and
\citet{semenova2022existence} establish that simpler models almost always
exist inside the Rashomon set alongside complex ones.

\paragraph{Lean formalization.}
The model selection impossibility is mechanically verified in two Lean~4
files.  \texttt{ModelSelection.lean} formalizes the \emph{Model Rashomon
Property}---for any two distinct candidate models, there exist evaluation
instances ranking them in opposite orders---and proves
\texttt{model\_selection\_impossibility}: no stable complete ranking
can be faithful to all evaluation instances.  The proof is a three-line
\texttt{linarith} argument from the Rashomon witnesses.
\texttt{ModelSelectionInstance.lean} instantiates the abstract
\texttt{ExplanationSystem} framework from
\texttt{ExplanationSystem.lean} and derives
\texttt{model\_selection\_instance\_impossibility} as a direct corollary
of the universal theorem via \texttt{explanation\_impossibility}.

\paragraph{Instance 6 (Model Selection Impossibility).}

\begin{theorem}[Model Selection Impossibility]
\label{thm:model-selection}
The model selection explanation system $\mathcal{S}_{\mathrm{sel}}$
(Definition~\ref{def:ms-system}) satisfies the Rashomon property
whenever the $\varepsilon$-Rashomon set has positive measure.
Consequently, by Theorem~\ref{thm:universal}, no selection rule
$\mathsf{exp}: \sH \to Y$ can simultaneously be:
\begin{enumerate}
  \item \textbf{Faithful:} $\mathsf{exp}(h) = \mathsf{obs}(h)$ for all $h$
    (the selected model is the one the evaluation metric actually endorses);
  \item \textbf{Stable:} $\mathsf{exp}(h) = \mathsf{exp}(h')$ whenever
    $\theta(h) = \theta(h')$ (the same configuration always yields the same
    selection); and
  \item \textbf{Decisive:} $\mathsf{exp}(h)$ is a strict recommendation
    for every $h$ (no ties or abstentions are reported).
\end{enumerate}
\end{theorem}

\paragraph{Design space.}
The model selection design space mirrors the attribution design space
exactly (proved in \texttt{ModelSelectionDesignSpace.lean}): every
selection rule falls into one of two families.
\begin{itemize}
  \item \textbf{Family~A} (decisive and unfaithful): the rule always
    picks a winner, but on some evaluation split it recommends the
    worse-performing model.  Standard single-criterion selection rules
    (e.g., argmax validation accuracy) fall here whenever the Rashomon
    set is non-trivial.
  \item \textbf{Family~B} (faithful and stable, not decisive): the rule
    declares a tie among Rashomon-set members, explicitly surfacing the
    ambiguity.  This corresponds to reporting the Rashomon set
    $\mathcal{R}_\varepsilon$ rather than a single selected model.
\end{itemize}

\paragraph{Implications.}
Theorem~\ref{thm:model-selection} formalizes and sharpens the empirical
observation of \citet{fisher2019models} that ``all models are wrong, but
many are useful'': when many models are equally useful by the available
metric, any decisive selection is \emph{necessarily} unfaithful on some
evaluation.  The instability practitioners observe when re-running
AutoML pipelines, changing random seeds, or switching validation folds
is not a software bug---it is the Rashomon property in action.

The resolution is the same as for feature attribution
(Section~\ref{sec:instance-attribution}): replace decisive selection
with a $G$-invariant aggregation over the Rashomon set.  Concretely,
this means reporting ensemble predictions across $\mathcal{R}_\varepsilon$,
using partial orders that render near-optimal models incomparable, or
applying Rashomon-aware model selection criteria
\citep{semenova2022existence,herren2023statistical}.  Each of these
trades decisiveness for the pair (faithful, stable), which
Theorem~\ref{thm:universal} shows is the only achievable combination
under underspecification.
